% Generated draft: requires mathematical and source review.
\documentclass[11pt]{article}
\usepackage[margin=25mm]{geometry}
\usepackage{fontspec}
\setmainfont{Latin Modern Roman}
\usepackage{amsmath,amssymb,mathtools,unicode-math}
\setmathfont{Latin Modern Math}
\usepackage{enumitem,needspace,xurl,hyperref}
\setlength{\parindent}{0pt}
\setlength{\parskip}{5pt}
\setlength{\emergencystretch}{3em}
\newcommand{\N}{\mathbb{N}}
\newcommand{\Z}{\mathbb{Z}}
\newcommand{\Q}{\mathbb{Q}}
\newcommand{\R}{\mathbb{R}}
\newcommand{\C}{\mathbb{C}}
\newcommand{\F}{\mathbb{F}}
\newcommand{\A}{\mathbb{A}}
\newcommand{\PP}{\mathbb{P}}
\newcommand{\E}{\mathbb{E}}
\newcommand{\eps}{\varepsilon}
\newcommand{\dd}{\,\mathrm{d}}
\begin{document}
\section*{UnsolvedMath ID: 4000004}
\textbf{Source problem code:} catalogue entry
\subsection{Definitions and mathematical statement}
Let $(X,d)$ be the intended metric state space and let $(m_x^t)_{x\in X,t\ge0}$ be transition laws for a continuous-time Markov semigroup, with $m_x^0=\delta_x$. When $m_x^t$ and $m_y^t$ are Borel probability measures with finite first moment, write
\[
 T_1(\mu,\nu)=\inf_{\pi\in\Pi(\mu,\nu)}
             \int_{X\times X}d(u,v)\,\mathrm d\pi(u,v),
\]
where $\Pi(\mu,\nu)$ is their set of couplings. For distinct $x,y\in X$, the proposed continuous-time coarse Ricci curvature is
\[
 \kappa(x,y)=\liminf_{t\to0^+}\frac1t
   \frac{d(x,y)-T_1(m_x^t,m_y^t)}{d(x,y)}.
\]
The state-space regularity and moment assumptions needed to make this expression and the process construction valid are not fully specified in the source.

Ollivier's Problem D asks whether this definition gives the standard elementary consequences of positive coarse Ricci curvature, under a natural assumption such as non-explosion, for both diffusion and jump processes, including processes with an unbounded generator. It also asks whether positivity of $\kappa$ alone implies non-explosion. For a process with lifetime $\zeta$, non-explosion means
\[
 \Pr_x(\zeta=\infty)=1\qquad\text{for every initial state }x.
\]
The source does not specify whether ``positive'' means $\kappa(x,y)>0$ for each distinct pair or a uniform bound $\kappa(x,y)\ge\kappa_0>0$. It does not enumerate the elementary consequences, give the class or domain of unbounded generators, or select the additional natural assumptions.

There is also a formulation issue in the second question: a conservative Markov semigroup already has probability transition laws, whereas a process killed at explosion generally yields subprobability kernels on $X$. The displayed transport distance is between probability laws. A non-explosion criterion starting from a possibly explosive generator therefore needs an explicit pre-explosion curvature definition, conservativity convention, or a specified cemetery-state construction. This dependency is left unresolved rather than assuming the desired conservativity in advance.
\subsection{Short English statement}
Does the infinitesimal transport definition of coarse Ricci curvature retain the usual positive-curvature consequences for continuous-time diffusion and jump processes with unbounded generators? Can curvature positivity itself prevent explosion?

The supplied research summary reports partial progress: the definition and many elementary consequences are established; positivity-to-non-explosion has partial results via Laplacian comparison; and much is known for jump processes. It says a fully general assertion for all unbounded generators is unresolved. No exact theorem, assumptions, or later primary reference is supplied for those progress claims. The 2008 problem and registry status are retained as supplied assessments, and provenance records no later authoritative resolution check.
\subsection{Sources}
\begin{itemize}
\item Local UnsolvedMath registry, permanent ID 4000004, source code catalogue entry, ``Continuous-time,'' proposed by Yann Ollivier, 2008. Dataset curation and metadata: Ulam AI and UnsolvedMath Contributors, CC BY 4.0.
\item Yann Ollivier, \emph{Discrete Ricci curvature: Open problems} (2008), Problem D, \url{https://www.yann-ollivier.org/rech/publs/problems_curvmarkov.pdf}; pdf-text extraction, accessed 2026-07-29.
\item The research summary supplies progress descriptions involving Laplacian comparison and jump processes, without separate bibliographic references or precise theorem statements. They were not independently checked here.
\item Source provenance marks status evidence NEEDS\_REVIEW and rights NEEDS\_REVIEW: the source presents the item as open, no later authoritative resolution check is recorded, and public access does not establish redistribution permission. This qualification is separate from the dataset licence.
\end{itemize}
\end{document}
